


Generated by \texttt{python -m crashguard.r7b\_decide} from \texttt{results/tables/r7b\_*.json} and \texttt{results/results.json} (keys \texttt{r7b.*}). Hypotheses as written in \texttt{docs/02\_preregistration/PREREGISTRATION\_V19.md}. The reserve size k = 4 was chosen from round-7 data (V19 section 0): this is a confirmation on fresh seeds, not a blind test.


\subsection*{1. Verdicts (Holm over two hypotheses, alpha = 0.05)}

{\scriptsize\raggedright
\begin{longtable}{p{0.190\linewidth}p{0.190\linewidth}p{0.190\linewidth}p{0.190\linewidth}p{0.190\linewidth}}
\toprule
Hypothesis & Criterion & Result & p & Verdict \\
\midrule
\endhead
H-R7b-1 (SIMULATED) & k = 4 policy: follow-up on time in at least 99\% of runs, lower 95\% limit, every feasible cell (9 cells, 1,200 runs each, seeds 100 to 105) & worst lower limit 99.8\%; rate at least 100.0\% in every cell & 5.78e-06 & SUPPORTED \\
H-R7b-2 (TRACE\_DRIVEN) & k = 4 policy against old, no attack: delivery -1 point at most and median time +10\% at most, every cell (18 cells, trace seed 4242) & delivery fails in 2 of 18 cells, time in 0 & 1 & NOT SUPPORTED \\
\bottomrule
\end{longtable}}

\subsection*{2. H-R7b-1 cells (fresh seeds, q = 1, primary stress)}

{\scriptsize\raggedright
\begin{longtable}{p{0.095\linewidth}p{0.095\linewidth}p{0.095\linewidth}p{0.095\linewidth}p{0.095\linewidth}p{0.095\linewidth}p{0.095\linewidth}p{0.095\linewidth}p{0.095\linewidth}p{0.095\linewidth}}
\toprule
Duty model & SF & n earlier events & old & k = 1 & k = 2 & k = 3 & k = 4 & k = 4 lower 95\% limit & median delay k = 4 (s) \\
\midrule
\endhead
H1 & 7 & 42 & 2.0\% & 72.8\% & 92.2\% & 99.8\% & 100.0\% & 99.8\% & 0.000 \\
H1 & 8 & 24 & 1.0\% & 29.7\% & 58.9\% & 89.0\% & 100.0\% & 99.8\% & 0.000 \\
H1 & 9 & 13 & 0.2\% & 36.0\% & 76.8\% & 100.0\% & 100.0\% & 99.8\% & 0.000 \\
H10 & 7 & 418 & 86.0\% & 100.0\% & 100.0\% & 100.0\% & 100.0\% & 99.8\% & 0.000 \\
H10 & 8 & 237 & 57.8\% & 100.0\% & 100.0\% & 100.0\% & 100.0\% & 99.8\% & 0.000 \\
H10 & 9 & 129 & 23.8\% & 100.0\% & 100.0\% & 100.0\% & 100.0\% & 99.8\% & 0.000 \\
T10 & 7 & 0 & 0.0\% & 100.0\% & 100.0\% & 100.0\% & 100.0\% & 99.8\% & 0.000 \\
T10 & 8 & 0 & 0.0\% & 100.0\% & 100.0\% & 100.0\% & 100.0\% & 99.8\% & 0.000 \\
T10 & 9 & 0 & 0.0\% & 100.0\% & 100.0\% & 100.0\% & 100.0\% & 99.8\% & 0.000 \\
\bottomrule
\end{longtable}}

The k = 1 column reproduces the V18 failure at the 1\% allowance (H1: 70\%, 27\%, 35\% in round 7) on fresh seeds. India's 10\% rules are the main deployment regime (H10 and T10 rows); 1\% (H1) is the stricter case. k = 2 and k = 3 are DESCRIPTIVE.

\subsection*{3. H-R7b-2: benign cost on the trace cells}

Cells that fail the criterion (delivery difference, k = 4 minus old, lower 95\% limit below -0.01):

{\scriptsize\raggedright
\begin{longtable}{p{0.136\linewidth}p{0.136\linewidth}p{0.136\linewidth}p{0.136\linewidth}p{0.136\linewidth}p{0.136\linewidth}p{0.136\linewidth}}
\toprule
Group & Scheme & Duty model & delivery difference & lower 95\% limit & median-time ratio & upper 95\% limit \\
\midrule
\endhead
H & R3 & T10 & -2.36 points & -2.42 points & 1.000 & 1.000 \\
H & R61 & T10 & -2.36 points & -2.42 points & 1.000 & 1.000 \\
\bottomrule
\end{longtable}}

**Trade-off of benign delivery against k (DESCRIPTIVE, group H, scheme R3, delivery difference against old in points; the other scheme and the other groups are in \texttt{results/tables/r7b\_trace\_cells.json}):**

{\scriptsize\raggedright
\begin{longtable}{p{0.190\linewidth}p{0.190\linewidth}p{0.190\linewidth}p{0.190\linewidth}p{0.190\linewidth}}
\toprule
Duty model & k = 1 & k = 2 & k = 3 & k = 4 \\
\midrule
\endhead
H1 & -0.40 & -0.40 & -0.40 & -0.40 \\
H10 & -0.40 & -0.40 & -0.40 & -0.40 \\
T10 & -2.36 & -2.36 & -2.36 & -2.36 \\
\bottomrule
\end{longtable}}

The curve is flat: the replay runs each start point as an isolated event with a fresh duty state, so the size of the reserve does not bind, and the cost comes from the withdrawal rule (a flash that cannot finish before the follow-up is due is not sent) and, under the off-time, from the time guard, which does not depend on k. A benign cost that grows with k would appear only under sustained load; that case was not run (NOT\_RUN).

\subsection*{4. Formal check (FORMAL, bounded)}

Bounds (V19 section 3): B in (40, 44, 48), W = 70, m in (1, 4), TC = 2, flash tries = 2, gaps (0, 2), earlier events 0 to 5, model T once per m.

{\scriptsize\raggedright
\begin{longtable}{p{0.317\linewidth}p{0.317\linewidth}p{0.317\linewidth}}
\toprule
Check & Statement & Result \\
\midrule
\endhead
F1 & Invariant 1, no class-L frame leaves less than R & holds, 0 violations \\
F2 & Invariant 2, reserve intact at pending, up to 5 earlier events & holds, 0 violations \\
F3 & k = 4 policy, unconditional Invariant 2, at most 3 earlier events & holds, 0 violations \\
F4 as written & k = 1 policy violates unconditional Invariant 2 with at most 3 earlier events & NOT met: 0 violations with at most 3 earlier events \\
F4 observed & at 4 earlier events and B = 40: k = 1 violates, k = 4 does not & yes \\
\bottomrule
\end{longtable}}

Per B (unconditional violations by the number of earlier events n; counts are weighted by the number of attacker choices that reach each state):

{\scriptsize\raggedright
\begin{longtable}{p{0.136\linewidth}p{0.136\linewidth}p{0.136\linewidth}p{0.136\linewidth}p{0.136\linewidth}p{0.136\linewidth}p{0.136\linewidth}}
\toprule
Model & B & policy & n = 3 & n = 4 & n = 5 & smallest n with a violation \\
\midrule
\endhead
H & 40 & old & 0 & 1963520 & 279425024 & 4 \\
H & 40 & k = 1 & 0 & 983040 & 172212224 & 4 \\
H & 40 & k = 4 & 0 & 0 & 45875200 & 5 \\
H & 44 & old & 0 & 176128 & 103211008 & 4 \\
H & 44 & k = 1 & 0 & 0 & 26157056 & 5 \\
H & 44 & k = 4 & 0 & 0 & 2490368 & 5 \\
H & 48 & old & 0 & 0 & 1343488 & 5 \\
H & 48 & k = 1 & 0 & 0 & 212992 & 5 \\
H & 48 & k = 4 & 0 & 0 & 0 & none \\
T & - & k = 1 & 0 & 0 & 0 & none \\
T & - & k = 4 & 0 & 0 & 0 & none \\
T & - & old & 655360 & 20971520 & 671088640 & 0 \\
\bottomrule
\end{longtable}}

Reading. The k = 4 policy keeps the follow-up on time through four earlier events at B = 40 and through five at B = 48; the k = 1 policy fails at four earlier events at B = 40. The old policy fails under the off-time from the start (T). Limits of the abstraction: the window of 70 steps holds about four events, so the old policy and the k = 1 policy do not fail before four earlier events, which is why control F4 as written (at most 3 earlier events) is not met; the distinction between the policies appears at the bound of four. The smallest violating n of the k = 4 policy is at 5 earlier events (V19 expected 4). The check is for the model, not for the radio.

